\begin{center}
    {\Large\bfseries Занятие 1. 4 класс}
\end{center}

\small
\begin{enumerate}[itemsep=0.7em]

\item \textbf{Разминка. Устный счёт «Змейка»}

Выполни действия по цепочке и запиши итоговый результат:

\begin{center}
\footnotesize
$15 \rightarrow [\times 4] \rightarrow [+24] \rightarrow [:3]
\rightarrow [+17] \rightarrow [-19] \rightarrow [\times 3]
\rightarrow [:2] \rightarrow [-14] \rightarrow \textbf{?}$
\end{center}

\item \textbf{Примеры на счёт}

В каждом примере необходимо выполнить ровно 3 действия. Обязательно соблюдай правильный порядок!

\begin{enumerate}[itemsep=0.35em]
\item Сначала сложение, затем умножение, затем вычитание:
\[
(2{,}345 + 1{,}655) \times 42 - 1{,}200 =
\]
\item Сначала умножение, затем вычитание, затем сложение:
\[
5{,}124 \times 36 - 14{,}464 + 8{,}000 =
\]
\item Сначала умножение, затем деление на круглое число, затем сложение:
\[
1{,}250 \times 48 : 60 + 900 =
\]
\end{enumerate}

\item \textbf{Задача Васи, Маши и Пети}

Васе задали на уроке 4 примера. У Маши он взял сами примеры, а Петя продиктовал ему ответы. Вот только Петя и Маша записали их в разном порядке! Но Вася догадался, как, совершенно не считая в столбик, правильно сопоставить примеры и ответы. Догадайтесь и вы! Соедините линией пример и его ответ.

\newpage

\vspace{0.25cm}
\noindent
\begin{minipage}[t]{0.47\linewidth}
\textbf{Примеры:}
\begin{itemize}[leftmargin=*,itemsep=0.15em]
\item[А)] $8{,}143{,}512 : 2$
\item[Б)] $3{,}136 \times 400$
\item[В)] $94{,}827 + 53{,}529$
\item[Г)] $48 \times 26{,}137$
\end{itemize}
\end{minipage}
\hfill
\begin{minipage}[t]{0.47\linewidth}
\textbf{Ответы (перепутаны):}
\begin{itemize}[leftmargin=*,itemsep=0.15em]
\item[1)] $1{,}254{,}400$
\item[2)] $4{,}071{,}756$
\item[3)] $1{,}254{,}576$
\item[4)] $148{,}356$
\end{itemize}
\end{minipage}

\item \textbf{Задачи на счёт}
\begin{enumerate}[itemsep=0.35em]
\item \dots
\item \dots
\end{enumerate}

\item \textbf{РАБОТАЕМ В ПАРАХ}
\begin{enumerate}[itemsep=0.5em]
\item Поставь между некоторыми цифрами знаки действий и скобки так, чтобы равенства стали верными.
\[
9 \quad 9 \quad 9 \quad 9 = 10
\]
\[
9 \quad 9 \quad 9 \quad 9 \quad 9 = 10
\]

\item Дано верное равенство:
\[
1 + 2 + 3 + 4 + 5 + 6 + 7 + 8 + 9 = 45
\]
Замени только один из знаков сложения на знак умножения так, чтобы значение выражения стало равным 100.
\end{enumerate}

\end{enumerate}
